\documentclass[11pt]{article}
\usepackage[margin=1in]{geometry}
\usepackage{amsmath,amssymb,amsthm}
\usepackage{hyperref}
\newtheorem{theorem}{Theorem}
\newtheorem{lemma}[theorem]{Lemma}
\theoremstyle{remark}
\newtheorem{remark}[theorem]{Remark}

\title{Proof of Conjecture 00000000035:\\
monochromatic Schur triples with $xy+1$ prime, $N=17$}
\author{jilint777}
\date{2026-10-04}

\begin{document}
\sloppy
\maketitle

\begin{abstract}
Conjecture 00000000035 asserts that there is an $N$ such that every
2-coloring of $[N]=\{1,\dots,N\}$ contains a monochromatic Schur triple
$(x,y,z)$, $x+y=z$, with $xy+1$ prime. We prove this with $N=17$, even when
$x<y$ is required. The value $17$ is optimal for $x<y$: an explicit coloring
of $[16]$ avoids all such triples. (If $x=y$ is allowed, $N=7$ already
suffices.) The proof is a finite case analysis, fully checked in Lean~4 (core
library only) and by exhaustive search in Python.
\end{abstract}

\section{Statement and reading}
The conjecture (English) reads:
\begin{quote}
There exists $N(k)$ such that any 2-coloring of $[N]$ contains a monochromatic
Schur triple $(x,y,z)$ with $x+y=z$ and $xy+1$ prime.
\end{quote}
The Chinese version is identical. A Schur triple in $[N]$ is
$x,y,z\in[N]$ with $x+y=z$. The parameter $k$ in ``$N(k)$'' does not occur
anywhere else in the statement, and the coloring is explicitly a
2-coloring. So the statement is: \emph{there exists $N$ such that every
$c:[N]\to\{R,B\}$ has $x,y,z\in[N]$ with $x+y=z$, $xy+1$ prime and
$c(x)=c(y)=c(z)$.} Since $N=17$ works for every value of $k$, any
dependence on $k$ is harmless. We prove the stronger version that requires
$x<y$, which in particular makes $x,y,z$ distinct.

\section{The triples}
Let $\mathcal T$ be the set of $(x,y,z)$ with $1\le x<y$, $z=x+y\le17$ and
$xy+1$ prime. There are exactly $31$ of them:
\begin{center}\small
$(1,2,3),(1,4,5),(1,6,7),(1,10,11),(1,12,13),(1,16,17),(2,3,5),(2,5,7),(2,6,8),$\\
$(2,8,10),(2,9,11),(2,11,13),(2,14,16),(2,15,17),(3,4,7),(3,6,9),(3,10,13),$\\
$(3,12,15),(3,14,17),(4,7,11),(4,9,13),(4,10,14),(4,13,17),(5,6,11),(5,8,13),$\\
$(5,12,17),(6,7,13),(6,10,16),(6,11,17),(7,10,17),(8,9,17)$.
\end{center}
For example $1\cdot2+1=3$, $8\cdot9+1=73$ and $7\cdot10+1=71$ are prime. Only the
\emph{presence} of these triples is used; that each $xy+1$ is prime is checked
in Lean with a verified trial-division checker.

\section{\texorpdfstring{Proof that $N=17$ works}{Proof that N=17 works}}
\begin{theorem}\label{thm:main}
Every 2-coloring $c$ of $[17]$ has a monochromatic triple in $\mathcal T$.
\end{theorem}
\begin{proof}
Suppose $c$ has none. Then for each $(x,y,z)\in\mathcal T$ the three colors
are not all equal. Hence, \emph{if two of them are equal, the third has the
other color}. Write ``$(x,y,z)\Rightarrow c(w)=\cdot$'' for this deduction.
We branch on a few values (a branching value is written without an arrow) and propagate. Every branch ends with a
triple of $\mathcal T$ whose three colors are equal, which is a contradiction. The
$12$ branches below cover all cases: at each bold split both colors are
taken (all 12 combinations that arise are listed).
\begin{enumerate}\small
\item \textbf{$c(2)=R$}, \textbf{$c(17)=R$}, $(2,15,17)\Rightarrow c(15)=B$, \textbf{$c(3)=R$}, $(1,2,3)\Rightarrow c(1)=B$, $(2,3,5)\Rightarrow c(5)=B$, $(1,4,5)\Rightarrow c(4)=R$, $(3,4,7)\Rightarrow c(7)=B$, $(1,6,7)\Rightarrow c(6)=R$, $(2,6,8)\Rightarrow c(8)=B$, $(3,6,9)\Rightarrow c(9)=B$, $(3,14,17)\Rightarrow c(14)=B$, $(4,13,17)\Rightarrow c(13)=B$, $(5,8,13)$ monochromatic.
\item \textbf{$c(2)=R$}, \textbf{$c(17)=R$}, $(2,15,17)\Rightarrow c(15)=B$, \textbf{$c(3)=B$}, $(3,12,15)\Rightarrow c(12)=R$, $(5,12,17)\Rightarrow c(5)=B$, \textbf{$c(6)=R$}, $(2,6,8)\Rightarrow c(8)=B$, $(5,8,13)\Rightarrow c(13)=R$, $(1,12,13)\Rightarrow c(1)=B$, $(1,4,5)\Rightarrow c(4)=R$, $(4,13,17)$ monochromatic.
\item \textbf{$c(2)=R$}, \textbf{$c(17)=R$}, $(2,15,17)\Rightarrow c(15)=B$, \textbf{$c(3)=B$}, $(3,12,15)\Rightarrow c(12)=R$, $(5,12,17)\Rightarrow c(5)=B$, \textbf{$c(6)=B$}, $(3,6,9)\Rightarrow c(9)=R$, $(2,9,11)\Rightarrow c(11)=B$, $(5,6,11)$ monochromatic.
\item \textbf{$c(2)=R$}, \textbf{$c(17)=B$}, \textbf{$c(3)=R$}, $(1,2,3)\Rightarrow c(1)=B$, $(1,16,17)\Rightarrow c(16)=R$, $(2,3,5)\Rightarrow c(5)=B$, $(1,4,5)\Rightarrow c(4)=R$, $(2,14,16)\Rightarrow c(14)=B$, $(3,4,7)\Rightarrow c(7)=B$, $(1,6,7)\Rightarrow c(6)=R$, $(2,6,8)\Rightarrow c(8)=B$, $(3,6,9)\Rightarrow c(9)=B$, $(8,9,17)$ monochromatic.
\item \textbf{$c(2)=R$}, \textbf{$c(17)=B$}, \textbf{$c(3)=B$}, $(3,14,17)\Rightarrow c(14)=R$, $(2,14,16)\Rightarrow c(16)=B$, $(1,16,17)\Rightarrow c(1)=R$, \textbf{$c(6)=R$}, $(1,6,7)\Rightarrow c(7)=B$, $(2,6,8)\Rightarrow c(8)=B$, $(3,4,7)\Rightarrow c(4)=R$, $(1,4,5)\Rightarrow c(5)=B$, $(4,10,14)\Rightarrow c(10)=B$, $(7,10,17)$ monochromatic.
\item \textbf{$c(2)=R$}, \textbf{$c(17)=B$}, \textbf{$c(3)=B$}, $(3,14,17)\Rightarrow c(14)=R$, $(2,14,16)\Rightarrow c(16)=B$, $(1,16,17)\Rightarrow c(1)=R$, \textbf{$c(6)=B$}, $(3,6,9)\Rightarrow c(9)=R$, $(2,9,11)\Rightarrow c(11)=B$, $(6,11,17)$ monochromatic.
\item \textbf{$c(2)=B$}, \textbf{$c(17)=R$}, \textbf{$c(3)=R$}, $(3,14,17)\Rightarrow c(14)=B$, $(2,14,16)\Rightarrow c(16)=R$, $(1,16,17)\Rightarrow c(1)=B$, \textbf{$c(6)=R$}, $(3,6,9)\Rightarrow c(9)=B$, $(2,9,11)\Rightarrow c(11)=R$, $(6,11,17)$ monochromatic.
\item \textbf{$c(2)=B$}, \textbf{$c(17)=R$}, \textbf{$c(3)=R$}, $(3,14,17)\Rightarrow c(14)=B$, $(2,14,16)\Rightarrow c(16)=R$, $(1,16,17)\Rightarrow c(1)=B$, \textbf{$c(6)=B$}, $(1,6,7)\Rightarrow c(7)=R$, $(2,6,8)\Rightarrow c(8)=R$, $(3,4,7)\Rightarrow c(4)=B$, $(1,4,5)\Rightarrow c(5)=R$, $(4,10,14)\Rightarrow c(10)=R$, $(7,10,17)$ monochromatic.
\item \textbf{$c(2)=B$}, \textbf{$c(17)=R$}, \textbf{$c(3)=B$}, $(1,2,3)\Rightarrow c(1)=R$, $(1,16,17)\Rightarrow c(16)=B$, $(2,3,5)\Rightarrow c(5)=R$, $(1,4,5)\Rightarrow c(4)=B$, $(2,14,16)\Rightarrow c(14)=R$, $(3,4,7)\Rightarrow c(7)=R$, $(1,6,7)\Rightarrow c(6)=B$, $(2,6,8)\Rightarrow c(8)=R$, $(3,6,9)\Rightarrow c(9)=R$, $(8,9,17)$ monochromatic.
\item \textbf{$c(2)=B$}, \textbf{$c(17)=B$}, $(2,15,17)\Rightarrow c(15)=R$, \textbf{$c(3)=R$}, $(3,12,15)\Rightarrow c(12)=B$, $(5,12,17)\Rightarrow c(5)=R$, \textbf{$c(6)=R$}, $(3,6,9)\Rightarrow c(9)=B$, $(2,9,11)\Rightarrow c(11)=R$, $(5,6,11)$ monochromatic.
\item \textbf{$c(2)=B$}, \textbf{$c(17)=B$}, $(2,15,17)\Rightarrow c(15)=R$, \textbf{$c(3)=R$}, $(3,12,15)\Rightarrow c(12)=B$, $(5,12,17)\Rightarrow c(5)=R$, \textbf{$c(6)=B$}, $(2,6,8)\Rightarrow c(8)=R$, $(5,8,13)\Rightarrow c(13)=B$, $(1,12,13)\Rightarrow c(1)=R$, $(1,4,5)\Rightarrow c(4)=B$, $(4,13,17)$ monochromatic.
\item \textbf{$c(2)=B$}, \textbf{$c(17)=B$}, $(2,15,17)\Rightarrow c(15)=R$, \textbf{$c(3)=B$}, $(1,2,3)\Rightarrow c(1)=R$, $(2,3,5)\Rightarrow c(5)=R$, $(1,4,5)\Rightarrow c(4)=B$, $(3,4,7)\Rightarrow c(7)=R$, $(1,6,7)\Rightarrow c(6)=B$, $(2,6,8)\Rightarrow c(8)=R$, $(3,6,9)\Rightarrow c(9)=R$, $(3,14,17)\Rightarrow c(14)=R$, $(4,13,17)\Rightarrow c(13)=R$, $(5,8,13)$ monochromatic.
\end{enumerate}
Each branch is a complete chain of forced deductions from the listed
triples, so no coloring avoids $\mathcal T$.
\end{proof}

\begin{theorem}
For every $k$ there is $N(k)$ (namely $17$) such that every 2-coloring of
$[N]$ contains a monochromatic Schur triple $(x,y,z)$, $x+y=z$, with $xy+1$
prime. Conjecture 00000000035 is true.
\end{theorem}

\section{Sharpness}
\begin{theorem}
The coloring of $1,2,\dots,16$ given by $0010101110110101$ ($0=R$, $1=B$) has no
monochromatic triple $(x,y,x+y)$ with $1\le x<y$, $x+y\le16$ and $xy+1$ prime.
So $17$ is the least admissible $N$ when $x<y$ is required.
\end{theorem}
\begin{proof}
There are $23$ such triples in $[16]$, and each is checked directly (Lean
\texttt{good16\_check}, Python \texttt{verify.py}).
\end{proof}

\begin{remark}
If $x=y$ is allowed, the extra triples $(x,x,2x)$ with $x^2+1$ prime appear
($x=1,2,4,6,\dots$), and the least $N$ drops to $7$ (\texttt{verify.py}). The
conjecture only asks for existence, so either reading is proved by $N=17$.
\end{remark}

\section{Lean formalization}
The directory \texttt{lean4/} is a self-contained Lean~4.19.0 project (core
library only, no Mathlib).
\begin{itemize}
\item \texttt{Prime n} is the standard definition ($n\ge2$ and every divisor is
$1$ or $n$). \texttt{prime\_of\_check} and \texttt{check\_of\_prime} prove that a
Boolean trial-division checker decides it.
\item \texttt{SchurPrimeTriple N x y z} means $1\le x<y$, $x+y=z\le N$ and
\texttt{Prime (x*y+1)}. Colorings are functions $c:\mathbb N\to$ \texttt{Bool};
only the values on $[N]$ matter.
\item \texttt{schur\_17}: every $c$ has a monochromatic
\texttt{SchurPrimeTriple 17}. The proof is exactly the case analysis above,
generated mechanically: 11 case splits, and each propagation step uses one
triple via the lemmas \texttt{force\_x}, \texttt{force\_y}, \texttt{force\_z}
and \texttt{clash}.
\item \texttt{conjecture\_00000000035}:
$\forall k,\ \exists N,\ \forall c,\ \exists x\,y\,z,$ \texttt{SchurPrimeTriple N x y z}
$\wedge\ c(x)=c(y)=c(z)$. \texttt{conjecture\_00000000035\_le} is the variant that
only requires $x\le y$.
\item \texttt{sixteen\_avoids}: the coloring \texttt{good16} of $[16]$ avoids all
triples.
\end{itemize}
\texttt{lake build} succeeds. There is no \texttt{sorry}, no
\texttt{native\_decide}, and no added axioms. \texttt{\#print axioms} shows at
most \texttt{propext}, \texttt{Classical.choice} and \texttt{Quot.sound}.

\section{Reproduction}
\begin{verbatim}
cd lean4 && lake build     # Lean 4.19.0, core only
python3 verify.py          # brute force over all 2^17 colorings
pdflatex report.tex
\end{verbatim}

\end{document}
